\section{A Big Bang towards ancestral futures and the common world}
\label{sec:futures}

\begin{orientacao}{Desenvolvimento posterior --- composição e temporalidade}
O \textit{Big Bang} do título precisa ser definido como operação do
ensaio, caso não se trate do conceito cosmológico. Evite sugerir uma
criação a partir do nada que apagaria histórias, relações e
ancestralidades. Os sete espaços seguintes permanecem editáveis.
\end{orientacao}

\atoplanejado{I}{Retomar o bar transformado. Mostrar que compor o comum
requer trabalho continuado, a partir de dependências já existentes;
definir o sentido do \textit{Big Bang} neste percurso.}

\atoplanejado{II}{Introduzir \refobra{krenak2022}. Desenvolver a
ancestralidade como relação com possibilidades de existência e de
futuro, com atenção à situação e à voz de Krenak. Evitar uma noção
genérica de sabedoria ancestral.}

\atoplanejado{III}{Introduzir \refobra{martins2021}. Relacionar tempo
espiralar, oralitura e corpo-tela às formas do ensaio e à circulação
de sua estória; reconhecer as matrizes afro-diaspóricas específicas
dessa elaboração.}

\atoplanejado{IV}{Fazer os corpos retomarem gestos e palavras do início
sob novas condições. Desenvolver o encontro entre Latour e Martins
sem reduzir uma teoria do corpo à outra. Mostrar convergências,
diferenças e transformações produzidas pela leitura conjunta.}

\atoplanejado{V}{Voltar à conta, ao rio e às condições de respirar.
Especificar obrigações, formas de participação e meios de sustentar
a composição do comum. A proposta precisa alcançar práticas além
da mudança de vocabulário.}

\atoplanejado{VI}{Confrontar um limite da composição: conflito,
perda, assimetria ou entidade que não caiba nos acordos. Distinguir
mundo comum de homogeneização ontológica e de consenso obrigatório.}

\atoplanejado{VII}{Retomar a pergunta inicial e transformar seu modo
de enunciação. Concluir sem apagar a poeira ou encerrar a
responsabilidade pelas mediações. Indicar o que o ensaio tornou
possível e o que permanece como tarefa.}

% Inserir aqui a narrativa e a argumentacao da terceira parte.
